\chapter{End-of-Life Protocol}
\label{ch:eol_protocol}

To align with modern sustainable engineering principles (ISO 14001) and international directives on Waste Electrical and Electronic Equipment (WEEE Directive 2012/19/EU), a comprehensive End-of-Life (EOL) protocol was established for the electromagnetic brake joint assembly. This protocol outlines design for disassembly, material recovery streams, hazardous substance management, and circular economy strategies.

\section{Design for Disassembly (DfD) Architecture}
The joint assembly incorporates mechanical modularity to ensure rapid non-destructive dismantling at the end of operational life. Threaded fasteners are standardized across the subassemblies to avoid the unnecessary use of permanent adhesive bonds or ultrasonic welds:
\begin{enumerate}[leftmargin=*]
	\item Removal of external wiring harnesses and detachment of the low-side Darlington switching driver PCB.
	\item Unfastening the retaining bolts to release the mild steel yoke, friction plate, and pressure tabs.
	\item Sliding the solenoids out of the yokes to disassemble the electromagnet.
\end{enumerate}

\section{Material Stream Classification and Circular Recycling}
Table~\ref{tab:eol_materials} defines the end-of-life material recovery routes for all primary structural and electrical constituents of the joint assembly.

\begin{table}[H]
	\centering
	\caption{End-of-Life Material Recycling and Disposal Protocol}
	\label{tab:eol_materials}
	\resizebox{\linewidth}{!}{%
		\begin{tabular}{|l|l|c|l|}
			\hline
			\textbf{Component} & \textbf{Material Group} & \textbf{\% Mass} & \textbf{Recommended EOL Recovery Route} \\ \hline
			Electromagnet Coils & High-Purity Copper (\SI{0.28}{mm}) & 32\% & Pyrometallurgical / Hydrometallurgical Recycling \\ \hline
			Yoke Core \& Tabs & 1018 Mild Steel Stock & 48\% & Ferrous Scrap Metal Smelting \& Reprocessing \\ \hline
			Bobbins \& Housing & PETG / Nylon Polymer & 12\% & Mechanical Polymer Granulating \& Regrind \\ \hline
			Switching Driver & FR4 Board / Lead-Free Solder & 6\% & Specialized E-Waste Pyrolysis \& Metal Extraction \\ \hline
			Insulation / Fasteners & Kapton / Stainless Steel & 2\% & Material Recovery / Thermal Energy Recovery \\ \hline
		\end{tabular}%
	}
\end{table}

\section{Hazardous Substance Mitigation and Environmental Safety}
To comply with Restriction of Hazardous Substances (RoHS Directive 2011/65/EU) guidelines, all electronic components (TIP41C power transistors, 2N3904 BJTs, and 1N4007 diodes) are specified as lead-free and halogen-free.

\begin{itemize}[leftmargin=*]
	\item The control electronics PCB must be diverted from municipal landfill streams and processed through accredited e-waste facilities to recover trace semiconductor materials.
	\item Thermal insulation, such as Kapton tape and magnet wire enamel coatings must undergo controlled high-temperature incineration with gas scrubbing to eliminate toxic volatile organic compound (VOC) emissions during metal recovery.
	\item \textbf{Circular Material Efficiency:} The overall recyclable mass fraction of the joint mechanism exceeds $92\%$, supporting closed-loop manufacturing frameworks for industrial metrology hardware.
\end{itemize}